% One schema change, one axis, three checks measured against it: what each
% one can actually see over the change's life, and where the two that exist
% locally leave a gap that lines up exactly with the failure in the sibling
% figure ch16-two-places-a-change-can-fail.tex.
%
% Same idiom as figures/tikz/ch01-one-field.tex: x=1mm/y=1mm, a single-weight
% axis with ticks, event nodes anchored south above it, a span bracket under
% the stages a label needs to explain. One lane, not two, because there is
% one change here, not two compared side by side; the span bracket is simply
% used three times, at three depths, so the brackets that overlap in x do not
% collide in y. No cross mark appears anywhere in this figure: nothing on
% this axis is a step a walk cannot take, only three different fields of
% view over the same four stages, so ch07-the-walk.tex's mark for a blocked
% step has no work to do here.
%
% The third bracket is the one addition. Dashing already means, in
% ch07-the-walk.tex and ch15-the-order-decides.tex, a span nobody actually
% walks; this figure reuses that meaning for a check nobody has written.
% It is also left open at the right, with no closing tick, rather than drawn
% as a complete bracket: a closed bracket over stages 1 to 2 would say a
% check exists that covers exactly that span, and it does not. The check
% this label describes would need to see stage 4, the operations clients
% actually send, while it is standing at stage 1, and nothing between those
% two points is on this axis today. An open end says the span is unclosed,
% which is the whole claim; closing it, even with a different mark, would
% say the opposite.
%
% No quotation marks anywhere in this file, including this comment block,
% per the note at the head of ch07-the-walk.tex and repeated in
% ch15-the-order-decides.tex.
%
% Bare tikzpicture (house style): the figure environment, caption and label
% live at the call site in the section file.
\begin{tikzpicture}[
  x=1mm, y=1mm,
  every node/.append style={font=\sffamily\scriptsize},
  axis/.style={draw, line width=0.4pt,
    -{Straight Barb[length=1.4mm, width=1.6mm]}},
  tick/.style={draw, line width=0.4pt},
  event/.style={anchor=south, align=center, text width=28mm, inner sep=1pt},
  span/.style={draw, line width=0.4pt},
  spanlabel/.style={anchor=north, align=center, inner sep=2pt},
  gap/.style={draw, line width=0.4pt, dashed},
]

\draw[axis] (0,0) -- (146,0);
\foreach \x in {15,54,93,131}{\draw[tick] (\x,-1.5) -- (\x,1.5);}

\node[event] at (15,3)  {you edit your\\subgraph document};
\node[event] at (54,3)  {composition merges\\four documents};
\node[event] at (93,3)  {the router loads\\the config};
\node[event, text width=32mm] at (131,3)
  {a client sends an\\operation it wrote\\last year};

% ------------------------------------------------- the two checks that exist
\draw[span] (47,-4) -- (47,-7) -- (61,-7) -- (61,-4);
\node[spanlabel, text width=40mm] at (54,-8)
  {composition. Reads the four documents and nothing else, so every
   question it can ask is a question about them};

\draw[span] (124,-4) -- (124,-7) -- (138,-7) -- (138,-4);
\node[spanlabel, text width=40mm] at (131,-8)
  {the router validates the operation. Correct, and after the
   deployment, in somebody else's request};

% ------------------------------------------- the check nobody has written
% Dashed throughout, and open at the right: no closing tick at stage 2,
% because the span this check would need does not close on this axis.
\draw[gap] (15,-22) -- (15,-25);
\draw[gap] (15,-25) -- (54,-25);
\node[spanlabel, text width=48mm] at (34,-26)
  {the check that would have caught it. Needs the operation on the right of
   this axis while standing at the edit on the left, and nothing in this book
   has it};

\end{tikzpicture}
